% 第6章 合成判据与操作流程
\section{四道核验：从判据到流程}\label{sec:criteria}

前两章各给了一条结构性命题，这一章把它们与成本、统计纪律合成一份可以照着执行的核验流程，并说清它与实务中“三道门”的对应关系。

\subsection{观测缝隙的定义}

\begin{definition}[观测缝隙]\label{def:oggap}
给定观测装置 $\mathcal{O}=(\mathbb{T},\mathcal{F})$（时间尺度 $\mathbb{T}=(\Delta_{\min},\tau_e,\nu_0)$，参照系 $\mathcal{F}$），称策略 $\pi$ 在该装置下存在\emph{观测缝隙}，当且仅当同时成立：
\begin{enumerate}
  \item \textbf{尺度可行}：存在网格 $\Delta\ge\Delta_{\min}$，使检测门 \eqref{eq:detectgate} 与盈利门 \eqref{eq:profitgate} 同时满足（即尺度窗口 \eqref{eq:window} 非空）；
  \item \textbf{参照系后仍为正}：$\alpha_{\mathcal{F}}(\pi)$ 在扣除成本 $C(\pi)$ 后为正，且超过检验阈值（含选择偏差校正）；
  \item \textbf{跨装置不反号}：在窗口内的尺度与参照系构成的网格上，净方向的符号不出现翻转。
\end{enumerate}
\end{definition}

需要如实说明这个定义的地位：它是合成式的，不是定理——三条都是可分别检验的必要条件，合起来在检验功效允许的范围内构成充分条件。“不反号”一条与统计功效纠缠（样本少时“没检出反号”不代表“不反号”），实际使用时应按 \ref{sec:pipeline} 小节的分级流程给出结论，而不是二元判定。合成的价值在于：\textbf{它把“这条信号有没有缝隙”从一个笼统的问句，拆成三个各有着落、可以分别拿数据回答的子问句。}

\subsection{四道核验}

把定义 \ref{def:oggap} 翻译成操作语言，就是四道核验。实务中常说的“三道门”（做空可实施性、换手成本、市场暴露）是其中两道半的合并版本；本文的版本单独列出尺度，并把统计校正从“隐含假设”提为独立一道。

\begin{table}[htbp]
\centering
\caption{理论条件与操作核验的映射}
\label{tab:gates}
\begin{tabular}{@{}>{\raggedright\arraybackslash}p{2.5cm}>{\raggedright\arraybackslash}p{5.6cm}>{\raggedright\arraybackslash}p{6.0cm}@{}}
\toprule
核验 & 做什么 & 判废信号（任一触发即降级） \\
\midrule
第一道：尺度 & 固定信号与成本口径，在网格序列 $\{1,2,3,5,10,20\}$（日）上重估净收益，画零成本、成本、净收益三条曲线 & 窗口为空（全部网格净收益 $\le 0$）；或峰位出现在窗口端点之外（说明模型或口径有误） \\[0.3em]
第二道：成本与实施 & 三档成本（零、费率、含冲击）分别记账；换手按真实调仓计算；逐股 ADV 与参与率上限；做空侧按市场制度判定可行性 & 含冲击净收益 $\le 0$；参与率超限的成交占比不可忽略；做空不可实施而未改用仅多头口径 \\[0.3em]
第三道：参照系 & 构造分解链（等权 $\to$ 市值加权 $\to$ 加行业/风格 $\to$ 加动量），记录 $\alpha$ 路径与其成本后的净值 & 在某件工具加入时 $\alpha$ 骤然归零（是暴露而非缝隙）；路径出现反号 \\[0.3em]
第四道：统计 & 试验账本（含信号候选、方向选择、网格选择、参数网格）；DSR 两口径；PBO 的 CSCV；$|t|>3$ 与 BH 校正 & 校正后不显著；PBO 接近或超过 $0.5$；试验账本无法说清 \\
\bottomrule
\end{tabular}
\end{table}

第三道里的分解链工具取舍，应当遵一个原则：\textbf{只用可实施、可交易的参照工具}。用不可交易的学术因子做正交，等于用纸上工具抹掉真实收益；这类“净化”得到的 $\alpha$ 只是残差的名字，不是可以提取的量。

统计校正的工具箱在文献里已有成熟形式——现实检验与高门槛\cite{white2000,sullivan1999,harveylz2016}、缩水夏普\cite{bailey2014dsr}与回测过拟合概率\cite{bailey2016pbo}，此处不重复推导；跨样本的复制研究给出的图景更细致——多数因子在扩展样本上仍成立，但幅度与显著性普遍下调\cite{jensen2023}。此外有一条容易漏掉的记账纪律：\textbf{网格选择本身也是选择}。在 $6$ 个尺度上挑选净收益最好的那个，与在 $61$ 个候选里挑最好的一样，会引入向上的选择偏差——第 \ref{sec:measurement} 章的尺度扫描若被用来“选一个最划算的频率”，它就必须计入试验账本，喂给 DSR 与 PBO 的试验数 $N$ 要相应放大。观测装置的自由度与统计校正的账本，在这一处咬合。

\subsection{执行流程}\label{sec:pipeline}

把四道核验排成固定顺序，得到一份可以照着走的流程。顺序有讲究：先便宜的后昂贵的，先否定的后肯定的。

\begin{enumerate}
  \item \textbf{写清装置}。数据频率、执行延迟、成本参数（佣金、税、冲击模型、单笔资金与 ADV 之比）各自填入固定表格。装置不写清，“有没有缝隙”这个问题不成立。
  \item \textbf{过必要条件}。用事件窗法估 $\hat\delta_0$、用交易记录估 $\hat c$，检查置信下界 \eqref{eq:clb}。这一关不过，后面全部免谈——但这一关很便宜。
  \item \textbf{画尺度曲线}。在预注册的网格序列上重估净收益，确认窗口与峰位。峰位若对网格顺序敏感（如删除某档后峰位跳变），先怀疑口径而不是庆祝发现。
  \item \textbf{记成本三档}。含冲击一档是判决档（冲击模型与平方根律的形式见 Almgren 与 Chriss 2000\cite{almgrenchriss2000}、Almgren 等 2005\cite{almgren2005}）；零成本与费率档只用来分解“成本吃掉多少”与“信号本身多少”。
  \item \textbf{走参照系分解链}。记录 $\alpha$ 路径；骤然归零者按暴露结案，缓慢衰减者保留候审。
  \item \textbf{做统计校正}。试验账本先于统计量：候选数、方向数、网格数、参数数，逐项记账；然后 DSR 与 PBO 才算得有意义。
  \item \textbf{分级出结论}。四道全过：可交易候选，进入资金分配讨论；过前三道而统计不足：列为“条件性候选”（等待更长样本或更干净的检验）；其余：归档并写明卡在哪一道——归档记录本身就是下一次研究的先验。
\end{enumerate}

这套流程全部走完仍不保证缝隙为真——它保证的是另一件事：\textbf{凡是被它放行的候选，都经过了审美之外的检验；凡是被它拦下的候选，拦截理由可以被逐条复查。}在低信噪比的领域里，后一句话比前一句有用得多。

\subsection{本章要点}

观测缝隙的定义把存在性拆成三条可检验的必要条件；四道核验（尺度、成本与实施、参照系、统计）是它的操作版本，与实务“三道门”的差别集中在两处——尺度被单独列出、统计校正被单独立账。执行流程按“先便宜后昂贵、先否定后肯定”排序，并有一条容易忽略的记账纪律：网格选择本身要计入试验账本。下一章把这套流程完整地跑一遍。